\section{Discussion}
\label{sec:discussion}

\paragraph{Plan once, execute open-loop.}
The biggest qualitative jump in performance came not from any particular
parameter value but from the architectural decision to build the swing
plan \emph{once} per rally and then execute it open-loop. Re-solving
inverse kinematics on every $1/240$\,s step exposes the controller to
small differences between consecutive IK solutions; with a redundant
manipulator the solver can hop between neighbouring null-space branches,
which the position-controlled joints faithfully track, and the resulting
chatter wrecks the Cartesian path. Locking the plan reduces the problem
to a single one-shot setpoint sequence that the joint controller can
realise smoothly.

\paragraph{Symmetry exploitation.}
The mirror-IK trick of Section~\ref{sub:ik} is, with hindsight, an
elementary application of the fact that the two robots share a
kinematic chain and differ only in base orientation. Despite this it
was the single fix that lifted Player~1's contact rate from $12\%$ to
$86\%$. The lesson is that PyBullet's IK is a numerical black box: even
with a deterministic seed, sufficiently symmetric problems can land on
different branches because the solver perceives the world frame, not
the base-local frame, and the rest-pose bias acts in joint space. When
the analytic geometry is symmetric, it is cheaper to enforce the
symmetry by construction than to hope the numerical solver discovers it.

\paragraph{Tuning the model, not just the controller.}
Two of the eight swept constants --- the effective restitution
coefficient $e_{B}$ and the horizontal damping factor $\eta_{V}$ ---
appear in the \emph{prediction} step rather than the control step.
They are parameters of the world model used inside the planner. The
sweep is therefore best understood not as finding the optimal
controller for a known plant, but as jointly fitting the world model
and the controller. The fact that $e_{B}=0.87$ won (rather than the
nominal $0.93$ that the simulator advertises) and that $\eta_{V}=0.82$
emerged from scratch are concrete evidence that a small mismatch
between assumed and realised physics matters more than the precise
choice of strike-plane location.

\paragraph{Connection to course material.}
The controller is built almost entirely from primitives covered in
ECE 275:
\begin{itemize}
  \item \emph{Forward and inverse kinematics} of a serial 7-DOF
        manipulator, including the use of damped least-squares with a
        null-space bias to resolve redundancy.
  \item \emph{Trajectory planning} as a sequence of waypoints in
        Cartesian space, realised via joint-space interpolation under a
        velocity-limited position controller.
  \item \emph{Rigid-body collision modelling} with Newtonian
        restitution to predict the post-bounce state of the ball.
  \item \emph{Projectile dynamics} for prediction between contacts,
        solved in closed form via the quadratic roots
        of~\eqref{eq:bounce}.
  \item \emph{Workspace and joint-limit awareness}, enforced both as
        hard limits in IK and as soft clamps on the strike target.
\end{itemize}
Notably, we did \emph{not} use the inverse-dynamics
$\tau = \mathbf{M}\ddot{\jq}+\mathbf{V}+\mathbf{G}$ that featured in
Assignment~02. Because the simulator exposes only position-mode
control, dynamic compensation happens inside the simulator's joint
servo and cannot be tightened from the outside. The trade-off is that
fine-grained shaping of impact forces is not available; we instead
shape the impact \emph{geometrically} through windup--follow-through
timing.

\paragraph{Limitations and what is left on the table.}
The most obvious open issues are:
\begin{enumerate}
  \item \emph{Lateral coverage}: the strike plane is fixed in $x$;
        balls with large $|v_{y}|$ at bounce time still produce
        excursions in $y$ that approach the workspace bound
        ($\pm 0.50$\,m). A second-order correction --- tilting the
        strike plane to track $v_{y}$ --- would help.
  \item \emph{Recovery between swings in two-player mode}: between the
        end of one rally and the start of the next, the controller
        relaxes to the ready pose; for short, fast rallies the arm is
        still in motion when the next ball is served. A short
        ``reset'' interpolation tied to the
        \texttt{WAITING\_SERVE} state would improve consistency.
  \item \emph{Spin}: the ball model includes no spin shaping. Even
        small Magnus or backspin handling could let the controller hit
        deeper without sailing.
  \item \emph{Adversarial behaviour}: in two-player mode the controller
        treats every incoming ball symmetrically and never targets a
        weakness of the opponent. A simple model that biased the
        return $y$-coordinate away from the opponent's last
        successful intercept point would convert some of the
        even rallies into wins.
\end{enumerate}
A complementary residual-RL pipeline (PPO over $\delta\in[-1,1]^{7}$
added to the classical action) was prototyped to investigate
items 1--3, but the present submission is fully classical.
